%%%%%%%%%%%%%%%%%%%%%%%%%%%%%%%%%%%%%%%
% رزومه محمد عباسی
% قالب XeLaTeX
% نسخه ۱.۰
% ترجمه شده برای زبان فارسی
%%%%%%%%%%%%%%%%%%%%%%%%%%%%%%%%%%%%%%%

% قبل از documentclass پکیج‌های مشکل‌ساز را بارگذاری کنید
\documentclass[a4paper]{deedy-persian} % تغییرات لازم برای کاغذ A4
\usepackage{xcolor}
\usepackage{xepersian} % بسته پشتیبانی فارسی

\settextfont{XB Niloofar} % تغییر به فونت مناسب فارسی
\setlatintextfont{Times New Roman} % فونت متون لاتین
\setdigitfont{XB Niloofar} % فونت اعداد فارسی


\begin{document}
\hypersetup{%
    pdfborder = {0 0 0}
}
%----------------------------------------------------------------------------------------
%	بخش عنوان
%----------------------------------------------------------------------------------------

\lastupdated % چاپ متن آخرین به‌روزرسانی در گوشه بالا
\namesection{محمد}{عباسی}{\begin{RTL} % نام شما
وب‌سایت: \href{zilabs.ir}{\bf http://zilabs.ir} | شماره تماس: 00989106091633\\
\href{mailto:absi1377@hotmail.com}{\lr{absi1377@hotmail.com}} | محمد عباسی چهارده چریک (لینکدین) % اطلاعات تماس
 | گیت‌هاب: \lr{ma14ch}
\end{RTL}
}

\vspace{1em} % فاصله میان عنوان و ستون‌ها

%----------------------------------------------------------------------------------------
%	ستون‌ها
%----------------------------------------------------------------------------------------

%\noindent
\begin{minipage}[t]{0.62\textwidth} % ستون سمت چپ، ۳۲٪ عرض متن صفحه

%------------------------------------------------
% تجربه کاری
%------------------------------------------------
\setlength{\itemsep}{5pt} % فاصله بین آیتم‌ها
\setlength{\topsep}{10pt} % فاصله بین لیست و متن قبل و بعد

\textcolor{gray}{\huge {تجربه کاری}}
\vspace{\topsep} % فاصله عمودی اضافی

%% بخش فراگستر با هاشور
\textbf{فراگستر} | مهندس هوش مصنوعی \\
\vspace{\itemsep} 
فروردین ۱۴۰۴ - اکنون (هیبرید) | تهران، ایران
\begin{RTL}
\begin{tightitemize}
\item توسعه و بهینه‌سازی سیستم‌های \lr{RAG} برای فراخوانی اسناد فارسی با استفاده از \lr{LangChain} و \lr{LangGraph}
\item پیاده‌سازی، بهبود و دیپلوی سرویس‌های \lr{OCR}، \lr{Summarizer}، استخراج ویژگی و \lr{Letter Writer} در مقیاس سازمانی
\item طراحی و بهبود سرویس‌های \lr{Text-to-Speech} و \lr{Speech-to-Text} با فاین‌تیون مدل‌ها و بهینه‌سازی پرامپت‌ها
\item ادغام و استقرار مدل‌های زبانی (\lr{Qwen}، \lr{Aya}، \lr{Gemma}، \lr{LLaMA}) بر روی زیرساخت‌های \lr{Docker}، \lr{Ollama} و \lr{VLLM}
\item طراحی و پیاده‌سازی ماژول‌های \lr{AgentCore} و \lr{Agentic Systems} شامل \lr{Long-Term Memory} و گراف‌های تعاملی
\item مدیریت و بهینه‌سازی زیرساخت‌ها شامل بهبود شبکه \lr{Docker}، استقرار \lr{Kong API Gateway}، و مدیریت دیتابیس \lr{PostgreSQL} و \lr{SQL Server}
\item انجام فاین‌تیون مدل‌های زبان و صوت (مانا \lr{TTS}، \lr{ASR}) و طراحی پایپ‌لاین آموزش و ارزیابی
\item مستندسازی کامل سرویس‌ها، مدیریت \lr{Git Repository}ها و پاک‌سازی کدها و داکیومنت‌ها
\end{tightitemize}
\end{RTL}
\vspace{\topsep}

%% بخش دانشگاه مالک اشتر با هاشور
\textbf{دانشگاه مالک اشتر} | محقق \\
\vspace{\itemsep} 
اسفند ۱۴۰۳ - اکنون (پاره‌وقت) | تهران، ایران
\begin{RTL}
\begin{tightitemize}
\item پردازش تصاویر ماهواره‌ای (SAR/ISAR)
\end{tightitemize}
\end{RTL}
\vspace{\topsep} 

% بخش جدید اضافه شده
\textbf{دمیس} | کارشناس هوش مصنوعی \\
\vspace{\itemsep} 
بهمن ۱۴۰۳ – فروردین ۱۴۰۴ | اصفهان، ایران
\vspace{\itemsep}
\begin{RTL}
\begin{tightitemize}
\item پروژه تبدیل صوت به متن 
\item پروژه دستیار پزشک 
\end{tightitemize}
\end{RTL}
\vspace{\topsep}

\textbf{گروه صنعتی انتخاب} | کارشناس هوش مصنوعی \\
\vspace{\itemsep} % فاصله عمودی اضافی
آبان ۱۴۰۲ – فروردین ۱۴۰۳ | اصفهان، ایران، SnowaTec
\vspace{\itemsep} % فاصله عمودی اضافی
\begin{tightitemize}
\item \lr{RAG} برای راهنماهای دستورالعمل
\item سیستم پیشنهاددهنده غذا با استفاده از تلگرام‌بات (در گیت‌هاب موجود است)
\item استخراج ویژگی و شناسایی موجودیت‌های نامی با تنظیم دقیق مدل‌های زبانی
\item تولید متن و استخراج ویژگی با استفاده از مدل‌های زبان بزرگ
\item ساخت مدل‌های ترجمه با استفاده از ترنسفورمرها (در گیت‌هاب موجود است)
\item توسعه سیستم پیشنهاد دهنده غذا مبتنی بر مواد غذایی 
\item \lr{ASR}، \lr{TTS}، مدل‌های زبان بزرگ و هوش مصنوعی مولد با استفاده از \lr{NVIDIA\texttrademark NEMO}
\item جمع‌آوری و تحلیل نظرات کاربران محصولات دیجی‌کالا و آمازون (در گیت‌هاب موجود است)
\item ...
\end{tightitemize}
\vspace{\topsep} % فاصله عمودی اضافی

\textbf{گروه صنعتی انتخاب} | کارشناس پایش تکنولوژی \\
\vspace{\itemsep} % فاصله عمودی اضافی
شهریور ۱۴۰۱ – آبان ۱۴۰۲ | اصفهان، ایران، SnowaTec
\vspace{\itemsep} % فاصله عمودی اضافی
\begin{tightitemize}
\item نظارت بر مراکز هوش مصنوعی شرکت‌های مرتبط با لوازم خانگی (بوش، سامسونگ، ال‌جی و غیره)
\item بررسی ساختار پلتفرم هوش مصنوعی و خدمات آنلاین شرکت‌های برجسته لوازم خانگی
\item جستجوی پتنت
\item مقایسه لوازم خانگی هوشمند و روندهای آینده
\item بنچ مارک محصولات حوزه لوازم خانگی
\item پایش ابزار های هوش مصنوعی در حوزه سازمانی
\item ...
\end{tightitemize}
\vspace{\topsep} % فاصله عمودی اضافی

\end{minipage} % پایان ستون سمت چپ
\hfill
\begin{minipage}[t]{0.33\textwidth} % ستون سمت چپ، ۶۶٪ عرض متن صفحه
%------------------------------------------------
% مهارت ها
%------------------------------------------------
\textcolor{gray}{\huge{مهارت ها}}
\vspace{\itemsep}
\subsection*{برنامه‌نویسی}
پایتون و نوت‌بوک آی‌پایتون (در زمینه یادگیری ماشین و یادگیری عمیق) \\
R (مبتدی) \\
  جاوااسکریپت، HTML و CSS  (سابقاً) \\
مای‌اس‌کیوال (سابقاً) \\


\subsection{یادگیری ماشین}
رگرسیون \\
تحلیل خوشه‌بندی \\
دسته‌بندی\\
\sectionspace % فاصله سفید بین بخش‌ها

\subsection{یادگیری عمیق}
شبکه‌های عصبی \\
CNN \\
GNN \\
RNN (مبتدی) \\
GANs (مبتدی) \\ 
 رمزگذار و رمزگشا \\
 ترانسفورمرها (متبدی)\\
... \\

\sectionspace % فاصله سفید بین بخش‌ها

\subsection{نرم‌افزارها}
رپیدماینر \\
Cvat \\
\lr{Label Studio} \\
LMStudio \\
\LaTeX \\
... \\

\subsection{سیستم‌عامل‌ها}
لینوکس (حرفه‌ای) \\
ویندوز \\

\sectionspace % فاصله سفید بین بخش‌ها

\subsection{پلتفرم‌ها}
Anaconda \\
NVIDIA\texttrademark NEMO \\
گیت \\
گیت‌هاب \\
داکر \\
LibRecommender (سیستم‌های پیشنهاد‌دهنده) \\
Cornac (سیستم‌های پیشنهاد‌دهنده)\\

\sectionspace % فاصله سفید بین بخش‌ها

\subsection{نرم افزار های کاربردی}
آفیس ( اکسل، ورد، پاورپوینت ) \\
کانوا \\
آدوب فتوشاپ\\
ادوب آدیشن\\
آدوب پریمایر\\


\sectionspace % فاصله سفید بین بخش‌ها

\end{minipage} % پایان ستون سمت راست
\newpage

%------------------------------------------------
% صفحه دوم
%------------------------------------------------
%\noindent
\begin{minipage}[t]{0.62\textwidth} % ستون سمت راست، ۶۲٪ عرض متن صفحه
%------------------------------------------------
% ادامه بخش تجربه کاری
%------------------------------------------------


\textbf{ایران اسمارتک} | تولیدکننده محتوا و طراح وب‌سایت \\
\vspace{\itemsep} % فاصله عمودی اضافی
۱۴۰۰ – ۱۴۰۱ | اصفهان، ایران
\vspace{\itemsep} % فاصله عمودی اضافی
\begin{tightitemize}
\item طراحی عناصر و پوسترها با فتوشاپ
\item توسعه‌دهنده وب‌سایت با استفاده از وردپرس
\end{tightitemize}
\vspace{\topsep} % فاصله عمودی اضافی


\textbf{فردا سنجش سرمایه (\lr{FSSco.})} | طراح وب‌سایت \\
\vspace{\itemsep} 
۱۳۹۴ – ۱۳۹۵ | اصفهان، ایران
\vspace{\itemsep} % فاصله عمودی اضافی
\begin{tightitemize}
\item توسعه‌دهنده وب‌سایت با استفاده از وردپرس
\end{tightitemize}
\vspace{\topsep} % فاصله عمودی اضافی

%------------------------------------------------
% تحصیلات
%------------------------------------------------
\textcolor{gray}{\huge {تحصیلات}}
\vspace{\topsep} % فاصله عمودی اضافی

\textbf{دانشگاه اصفهان} | کارشناسی ارشد هوش مصنوعی  \\
\vspace{\itemsep} 
1401 – 1403 | اصفهان، ایران
\vspace{\itemsep} 
\begin{tightitemize}
\item الگوریتم‌های پیشرفته، نظریه محاسبات پیشرفته، هوش مصنوعی پیشرفته، یادگیری ماشین، مبانی نظری رمزنگاری، داده‌کاوی محاسباتی، یادگیری عمیق
\item سیستم پیشنهاددهنده غذا از با استفاده از داده‌کاوی عادت های تغذیه ای کاربران (پایان‌نامه)
\item چت‌بات صوت به صوت فارسی (موجود در گیت‌هاب)
\item سیستم پیشنهاددهنده کلان‌داده با استفاده از \lr{TwoTower}, \lr{DeepFM} و ...
\item سیستم پیشنهاد‌دهنده غذای سالم بر اساس دانش تغذیه‌ای
\item مروری بر سیستم‌های پیشنهاددهنده غذا
\item پیش‌پردازش کلان‌داده برای اهداف یادگیری عمیق و ماشین
\item ...
\end{tightitemize}
\sectionspace % فضای سفید بعد از بخش
\textbf{دانشگاه اصفهان} | کارشناسی علوم کامپیوتر  \\
\vspace{\itemsep} % فاصله عمودی اضافی
1396 – 1401 | اصفهان، ایران
\begin{tightitemize}
\item مبانی کامپیوتر و برنامه‌نویسی، برنامه‌نویسی پیشرفته، مبانی علوم ریاضی، ریاضیات عمومی 1 و 2، اصول سیستم‌های کامپیوتری، مبانی ترکیبیات، مبانی ماتریس‌ها و جبر خطی، مبانی منطق و نظریه مجموعه‌ها، مبانی تحلیل عددی، اصول سیستم‌عامل‌ها، طراحی و تحلیل الگوریتم‌ها، پایگاه داده‌ها، مبانی احتمال، معادلات دیفرانسیل، بهینه‌سازی خطی، نظریه محاسبات، مبانی جبر، رمزنگاری، نظریه کدگذاری، لاتک، متلب، موضوعات خاص در علوم کامپیوتر (هوش مصنوعی)، مدارهای منطقی
\end{tightitemize}
\sectionspace % فضای سفید بعد از بخش
\vspace{\topsep}
\textcolor{gray}{\huge {افتخارات و مقالات}}
\vspace{\itemsep} % فاصله عمودی اضافی
\begin{tightitemize}
\item ۱۵ درصد برتر ورودی هنگام فارغ التحصیلی در مقطع کارشناسی | دانشگاه اصفهان
\item سیستم پیشنهاددهنده غذای سالم از طریق داده‌کاوی عادت های تغذیه‌ای کاربران | ۱۵‌امین کنفرانس بین المللی فنآوری اطلاعات و دانش \lr{(IKT)}
\end{tightitemize}
\sectionspace % فضای سفید بعد از بخش
\vspace{\topsep}
\textcolor{gray}{\huge {دوره های گذرانده شده}}
\vspace{\itemsep} % فاصله عمودی اضافی
\begin{tightitemize}
\item گیت - GreatLearning (دارای گواهینامه)
\item هکاتون LLM - Quera (دارای گواهینامه)
\item مبانی علم داده و داده‌کاوی - Lynda
\item برنامه‌نویسی R از مبتدی تا پیشرفته همراه با تمرینات واقعی - Udemy
\item یادگیری برنامه‌نویسی پیشرفته پایتون 2020 - Udemy
\item CS50AI 2023 - دانشگاه هاروارد
\end{tightitemize}
\sectionspace % فضای سفید بعد از بخش
\vspace{\topsep}
\textcolor{gray}{\huge {کتاب های مطالعه شده}}
\vspace{\itemsep} % فاصله عمودی اضافی
\begin{tightitemize}
\item داده‌کاوی: مفاهیم و تکنیک‌ها (تألیف \lr{Jiawei Han})
\item مقدمه‌ای بر الگوریتم‌ها (CLRS)
\item هوش مصنوعی - رویکردی مدرن (2010, \lr{Prentice Hall})
\item سیستم‌های پشتیبانی تصمیم برای هوش تجاری، ویرایش دوم (2010, \lr{Wiley})
\item روش‌های ماتریسی در داده‌کاوی و شناسایی الگو (\lr{Lars Eldén})
\item یادگیری عمیق با پایتون (تألیف \lr{François Chollet})
\end{tightitemize}
\vspace{\topsep} 


\end{minipage} % پایان ستون سمت راست 
\hfill
\begin{minipage}[t]{0.33\textwidth} % ستون سمت چپ ۳۳٪ عرض متن صفحه
%------------------------------------------------
% علاقه مندی ها
%------------------------------------------------
\textcolor{gray}{\huge{علاقه مندی ها}}
\vspace{\itemsep}

\subsection*{علاقه‌مندی‌های پژوهشی}
مدل‌های زبانی بزرگ\\
سیستم‌های پیشنهاددهنده\\
پردازش تصویر SAR و ISAR\\
آینده‌پژوهی هوش مصنوعی\\
سیاست‌گذاری در هوش مصنوعی\\
صنعت \lr{4.0}\\

\sectionspace % فضای سفید بعد از بخش
\subsection*{فن‌آوری‌های اخیراً استفاده‌شده برای توسعه نرم‌افزار}
\textbf{زبان‌ها و چارچوب‌ها:}  
پایتون (فست‌ای‌پی‌آی، فلسک، استریم‌لیت)، جاوااسکریپت (نود.جی‌اس)، اچ‌تی‌ام‌ال/سی‌اس‌اس، بش اسکریپت  

\textbf{زیرساخت و استقرار:}  
داکر، داکر کامپوز، داکر سوآرم، اولاما، وی‌ال‌ال‌ام، لاما فکتوری، گیت‌وی کنگ، دِک، انجین‌اکس، آپاچی سرور، کلادفلر، اس‌کیوال سرور، پستگرس‌کیوال، پی‌جی‌ادمین  

\textbf{یادگیری ماشین و مدل‌های زبانی:}  
لنگ‌چین، لنگ‌گراف، هاگینگ‌فیس ترنسفورمرز، پای‌تورچ، تنسورفلو، انویدیا نمو، سنسنس ترنسفورمرز، ویسپر، کوئن، آایا، لاما، جِما، میسترال، جی‌پی‌تی ای‌پی‌آی، کوهِر  

\textbf{پردازش تصویر و متن:}  
اوپن‌سی‌وی، تسراکت او‌سی‌آر، پایپ‌لاین او‌سی‌آر، اکسترکشن ویژگی، سامرایزر، سرویس اعتبارسنجی، مترادف‌ساز  

\textbf{مدیریت داده و آموزش مدل:}  
پایپ‌لاین فاین‌تیون، پِفت، فلش اتنشن، پیش‌پردازش داده‌ها (ای‌اس‌آر، او‌سی‌آر، تی‌تی‌اس)، آناکوندا، پانداس، نام‌پای  

\textbf{مدیریت نسخه و مستندسازی:}  
گیت، گیت‌هاب، مارک‌داون، لاتک، استراکچر.ام‌دی، بهینه‌سازی ریدمی  

\textbf{ابزارهای توسعه و تست:}  
وی‌اس کد، پست‌من، اَپی اکسپلورر، گرادیو، فلاک، ال‌ام استودیو  

\sectionspace % فضای سفید بعد از بخش


\subsection*{مدل های زبانی (متنی) استفاده شده برای توسعه}
\lr{Qwen 2.5 \\
Qwen 3 \\
Cohere-aya \\
Llama 3 \\
mt5 \\
sentence-transform base models\\
}

\sectionspace % فضای سفید بعد از بخش
\subsection*{مدل های زبانی (متن و تصویر) استفاده شده برای توسعه}
\lr{Qwen 2.5vl \\
Mistral\\
}
\sectionspace % فضای سفید بعد از بخش

\subsection*{سرگرمی‌ها}
عکاسی\\
دوچرخه‌سواری\\
دویدن\\
تئاتر\\
سینما\\

\sectionspace % فضای سفید بعد از بخش

\subsection*{زبان‌ها}
فارسی (زبان مادری)\\
انگلیسی (مسلط)\\

\sectionspace % فاصله سفید بین بخش‌ها

\end{minipage} % پایان ستون سمت راست
\hrule

\end{document}
